% ---------------------------------------------------------------------------
% Introduction
%
% The chain, in order. Each link is one paragraph and each exists to make the
% next one necessary:
%
%   1. RL's home turf is adversarial games -- so the comparison happens where
%      RL is strongest, not where it is weakest.
%   2. Concede the extensibility objection, then relocate it: it describes the
%      author, not the artifact.
%   3. Coding agents change who the author is. The operative property is
%      editability, not interpretability.
%   4. The question, announced as two-sided.
%   5. Lower bound: information per sample. The on-policy/off-policy spectrum,
%      cross-policy-class data reuse, rules as a non-sample prior. The spectrum
%      is also the reason model-based RL has to be the primary baseline: against
%      policy gradient alone, symbolic-vs-neural is confounded with
%      on-policy-vs-off-policy. Keep that sentence attached to this link -- it
%      reads as an unmotivated methods note anywhere else.
%      This link ends by giving away more than it looks like: the
%      information-per-sample gain is NOT the symbolic form's doing, since a
%      memory-based agent holding no program gets it too (echo2025hindsight).
%      Conceding that is what keeps the two bounds independent, which is the
%      whole structure of the paper -- the lower bound is about what the loop
%      can read, the upper about what it can edit, and only the second is the
%      program's. Deleting the concession to strengthen the lower bound would
%      collapse them into one claim and hand a referee the objection ("then why
%      the program at all?") with no answer in the text.
%      Keep the system NAMES in both places -- "ECHO" here and "DiPRL" in the
%      editability paragraph. Chengpeng Hu (diprl2026) and Michael Y. Hu
%      (echo2025hindsight) are different people, and natbib renders both as
%      "Hu et al." with only the year to tell them apart, so the two read as one
%      group's pair. Naming the systems is what disambiguates them for a reader;
%      do not revert either to "a recent system". Found by reading the rendered
%      intro end to end -- no mechanical check catches it.
%   6. Upper bound, in four paragraphs: the mechanism (accumulation defeats
%      placement); the external evidence for it (SlopCodeBench); the two results
%      isolating the mechanisms underneath, which are about context and
%      execution rather than code quality (longcontextbugfix2026,
%      longhorizonexecution2026); and the distinction that keeps our loop
%      tractable (their spec grows, a game's rules do not, so a heuristic can be
%      replaced rather than layered). This was three paragraphs until the
%      citation check expanded the evidence; do not merge them back, the
%      combined version ran twenty rendered lines and carried three papers.
%      All three citations were verified against full text 2026-08-18 and all
%      three needed correction -- see their refs.bib annotations, which are
%      authoritative, and note in particular that longhorizonexecution2026's
%      own thesis runs opposite to the use we make of it.
%      Two claims about SlopCodeBench were drafted wrong and are now fixed
%      against the v2 full text; both flattered us, which is why they lasted.
%      (a) Human repos do NOT "flatten". 68% of them end more verbose than they
%      began. The finding is a rate gap -- per-checkpoint slope 0.0022 vs 0.0144
%      for verbosity -- and the rate gap is the better claim anyway, since a
%      human exemption would make maintenance capacity look like an
%      agent-specific defect instead of a general limit that agents hit sooner.
%      (b) Never write that quality-aware prompting leaves the slope
%      "statistically unchanged". Their paper runs no significance test on that
%      experiment at all; "slope" and "intercept" appear only in Appendix D,
%      about the human panel. Say the prompts move the starting level without
%      slowing the growth -- and keep the price attached (2.3 pp of strict
%      correctness, 12.1% cost), because an intervention that trades capability
%      for initial quality and still does not bend the curve is the actual
%      argument for a capacity limit. All figures live in the refs.bib
%      annotation, which is authoritative; re-verify if a v3 appears, since
%      every v1 number changed in v2 and this paragraph inherited the v1 set.
%      Then the axis: description length of a good-enough policy against the
%      agent's maintenance capacity. Go is the counterexample that kills rule
%      complexity as the x-axis. The LLM boundary falls out of the same axis.
%      The inequality must state its observable signature where it is introduced
%      -- convergent stall size across games, divergent Elo at that size -- and
%      must say that it is the PAIR that does the diagnosing. Until 2026-08-18
%      this paragraph called the plateau "diagnostic" without saying how it could
%      be observed, and the signature appeared only in the contributions, which
%      asserted the result with no mechanism attached. Convergent length alone
%      proves nothing: it is equally consistent with games that never differed in
%      complexity.
%   7. Why the question is still open. The heading said "Why this ceiling has not
%      been measured" until 2026-08-18, which is this link's own prohibition
%      stated at heading level, and which its first clause then contradicted.
%      Headings get written early and revised least: check them against these
%      notes separately from the body. Careful: agent self-improvement
%      decay HAS been measured (frontiereng2026 -- frequency t^-1 AND magnitude
%      k^-1, marginal returns near zero after ~50-100 iterations, tail to 500).
%      Do not compress this to "plateau within tens of iterations", which the
%      paragraph said until 2026-08-18: it is the wrong number and the wrong
%      shape. Cite both exponents. Do not write "never measured" -- concede the
%      result and say what it lacks: single-agent tasks, no human field, no RL
%      baseline. The
%      confound we remove is a different one, human iteration cost.
%   8. The venue.
%   9. HL, plus the selection-effect concession.
%  10. Contributions.
%
% Four things to hold onto while editing:
%   - The objection is conceded, never refuted. "Everyone was wrong" is both
%     weaker and less true than "everyone was measuring the author".
%   - Paragraphs 5 and 6 stay adjacent. Split them and the paper reads as two
%     unrelated findings instead of one bound with two sides.
%   - Do not write "if-else". Characterizing what the winning programs contain
%     is a result, not a punchline.
%   - Do not claim editability as our coinage. Programmatic RL already calls
%     its policies human-editable; what changes here is the role the property
%     plays -- precondition for the optimizer, not feature of the deliverable.
%   - In paragraph 6, never let rule complexity back in as the axis. An earlier
%     draft rejected it with the Go example and then wrote "rule complexity and
%     revision count are the same axis seen twice" three sentences later, which
%     contradicts the paragraph's own thesis. The axis is K(pi_good), always.
%     A game defeats the loop because its *winning policy* is long, not because
%     its rulebook is.
%     Related, and also drafted wrong once: do not write that a hard game
%     "starts near the capacity". The program grows from a naive draft in every
%     game; the capacity is fixed by the agent. What varies across games is where
%     K(pi_good) sits relative to that ceiling -- below it, the program reaches
%     winning strength before it becomes unmaintainable; above it, the program
%     stops improving while still short of winning. Getting this backwards makes
%     the paper's central inequality read as though the game moved the ceiling.
%
% LENGTH. Roughly 2,200 words plus 400 of contributions, which is long. It is
% long because the argument is two-sided, carries external evidence, and answers
% two objections before the results -- all deliberate. Do not trim it to a target
% without a venue actually demanding one. When one does, cut in this order:
%   1. The editability discussion in para 3, from four sentences to two. The
%      distinction is needed; the elaboration is not.
%   2. The contamination-resistance contrast, to one sentence. It is positioning
%      against a literature, and the benchmark section makes the same case.
%   3. The two objection paragraphs, merged and moved to Limitations. This costs
%      real credibility -- a reader who meets the pretraining objection first in
%      Section 6 has already decided we were hiding it -- so it goes last.
% Never cut: the Go example (it is what kills rule complexity as the axis), the
% single inequality, or the replacement-vs-accumulation distinction.
% ---------------------------------------------------------------------------

\section{Introduction}

Reinforcement learning (RL) established itself on adversarial games. Chess, Go,
shogi, StarCraft, and Dota were where learned policies first beat the strongest
humans alive, and where the field's characteristic result took shape: a policy
trained by self-play, given enough interaction, exceeds anything a person wrote by
hand \citep{alphazero2018,alphastar2019,openaifive2019}. Games are also where the
symbolic alternative --- rules, an evaluation function, a search procedure --- was
judged to have the lower ceiling, and this paper stays on that turf rather than
retreating to a setting chosen to favour programs.

The objection is that a program exhibits only the behaviour somebody wrote down:
every new case has to be noticed, encoded, and reconciled with the rules already
present, at a cost that grows with the program. This is correct and we do not
dispute it. But it is a claim about the author, not the artifact. Nothing in the
symbolic paradigm caps what a program can express --- the cap is on how much of it
a human will maintain before the reconciliation cost exceeds what they can pay.
Asymptotically the learned policy still wins. The question is what happens under a
budget, where the authoring cost is the binding term.

\paragraph{What coding agents changed.} They change who the author is, and the
property that matters once the author is a machine is not the one usually
argued. The axis is \emph{editability}, not interpretability: interpretability
asks whether a human understands the policy, which a reader may reasonably not
care about, while editability asks whether the optimization loop can act on the
artifact at all. A symbolic program can be read in full and changed in part, and
takes the heterogeneous inputs maintenance requires --- rules, its own source,
replays of losses. A learned policy admits no comparable loop: its behaviour lives
in weights that cannot be edited locally and offer no handle by which a specific
failure traces to a specific parameter. The two properties do come apart --- a
system can reason in natural language about every decision it makes and still hold
its policy in weights \citep{tig2025think} --- and programmatic RL already calls
its policies human-readable and editable \citep[DiPRL;][]{diprl2026}. What changes
here is the role the property plays: not a feature of the delivered artifact, but
the reason an optimizer exists for it.

So the question reopens, and this time it is measurable: \emph{how far does the
symbolic paradigm go once a machine pays the maintenance cost?} The answer is
two-sided --- an agent-maintained program wins for a reason having nothing to do
with expressiveness, and stops for a reason having everything to do with the program
and nothing to do with the paradigm.

\paragraph{Why an agent-maintained program wins under a budget.} The binding
constraint under a budget is information per sample. A scalar reward reports that
an episode went badly; the rules plus a replay of the loss report what happened,
when the position turned, and which decision was responsible --- and the
difference compounds, because it determines how much of the sample survives to the
next revision. Policy-gradient methods are on-policy, using a batch for one update
and discarding it \citep{offpolicyvaluellm2026}; off-policy methods reuse it
\citep{samplereuse2022}. A maintained program sits at the far end of that
spectrum: off-policy and \emph{cross-policy-class}, able to learn from replays of
any opponent and from human source code no gradient method can consume, and it
reads the rules directly --- a prior obtained without spending samples at all. An
LLM revising a Python policy from execution traces is already competitive with
deep-RL baselines on single-agent tasks at far fewer environment interactions
\citep{kuang2025generative}, though only directionally: one family of tasks, no
shared budget accounting, no adversary. The spectrum also settles how the
comparison must be run, since measuring against policy gradient alone confounds
symbolic-versus-neural with on-policy-versus-off-policy --- so model-based RL is
our primary baseline and policy gradient the on-policy reference point.

The advantage is narrower than it looks, and the narrow version is the one worth
defending. ECHO rewrites its own failed episodes into usable experience and becomes
more sample-efficient holding no program at all \citep{echo2025hindsight}: what
buys the information per sample is the reader, not the symbolic form. The program
earns its place on the other side of the question, as what makes the policy an
object the loop can edit --- and therefore as what sets where the loop stops. The
two bounds are measured against different things: the lower against RL's sample
budget, the upper against the agent's own capacity.

\paragraph{Why it stops.} What the loop accumulates is what eventually stops it.
Revisions leave the program larger and more interconnected, and each further
revision has to place an edit without disturbing what already works; past some size
that placement fails more often than it succeeds --- not because the paradigm ran
out of expressiveness, but because the artifact outgrew what its maintainer can
hold, attribute, and change. This is measurable outside our setting. Coding agents
extending their own prior
code accumulate structural damage as they go, at rising cost and no matching gain
in correctness; human maintainers are not exempt but accumulate it several times
more slowly per revision, so what separates a maintainer that keeps working from
one that stalls is a rate rather than an exemption
\citep{slopcodebench2026}. Quality-aware prompting lowers the starting level
without slowing the growth, and costs correctness to do it --- the signature of a
capacity limit rather than a prompting defect. Two mechanisms underneath are not
about code quality at all: agentic repair succeeds only while the context each step
accumulates stays small, which locates the advantage in decomposition rather than
retrieval \citep{longcontextbugfix2026}; and given the plan up front, per-step
accuracy still falls as steps accumulate, models erring more once their own earlier
errors are in context \citep{longhorizonexecution2026}. Reasoning mitigates that
second failure, which helps us rather than hurts --- a plateau reached by a
reasoning agent cannot be attributed to self-conditioning --- and what survives is
the finite single-turn execution horizon each frontier model has, an order of
magnitude apart between them. That is capacity as a property of the agent rather
than of the task, and a program under revision is the case that cannot be
decomposed away: the artifact is the accumulated context, and every revision has to
hold it whole.

The evidence cuts both ways, though. Their decay is driven by a specification that
keeps growing, so nothing may be discarded; a game's rules do not change, so a
better heuristic may \emph{replace} a worse one instead of being layered on top of
it and the program is free to get shorter. Accumulation is a tendency here rather
than an obligation --- which is why the ceiling is worth measuring instead of
assuming, and why we log program size at every revision rather than only Elo.

The axis this limit falls on is a description length, and not the length of the
rules. Go's rules fit on a page, yet no compact symbolic policy plays it well ---
the programs that do carry their strength in weights --- while a contest game may
run to pages of special cases and still admit a strong policy in a few hundred
lines. What governs the outcome is the description length of a \emph{symbolic}
policy good enough to win, which we write $K(\pi_{\mathrm{good}})$ after
Kolmogorov complexity and estimate through a compression proxy, the quantity itself
being uncomputable \citep{lividanyi2019kolmogorov,rissanen1978mdl}. The prediction
is a single inequality: HL succeeds while $K(\pi_{\mathrm{good}})$ stays below the
agent's maintenance capacity, and plateaus as it approaches it. Since the program
grows from a naive draft in every game, what differs is where the target sits
relative to the ceiling --- below it, the program reaches winning strength before it
becomes unmaintainable; above it, the program stops improving while still short.
And because the capacity belongs to the agent, this predicts a signature: the size
at which the loop stalls should be roughly the same from game to game while the
strength reached there should not, and it is that pairing, not either half, that
distinguishes a limit of the maintainer from a limit of the game. The paradigm's
boundary then follows without a separate argument: tasks at the scale of language
modelling lie outside it not because they are adversarial or unstructured, but
because no policy of a maintainable description length is any good at them. And
since the ceiling belongs to the maintainer, it moves with model capability, which
is why it can be measured at all.

\paragraph{Why the question is still open.} Agent self-improvement is already
known to decay, and from both directions at once: across 47 real engineering
tasks, improvement frequency falls as $t^{-1}$ while the magnitude of the $k$-th
improvement falls as $k^{-1}$, a double squeeze driving marginal returns near zero
after 50 to 100 iterations \citep{frontiereng2026}. What that leaves open is where
the symbolic paradigm stands against a learned one, since those tasks are
single-agent and carry neither a human field nor an RL baseline. Where both
reference points were present, the maintainer was always human, so the paradigm
and the cost of human iteration were measured together: a heuristic program that
stopped improving might have exhausted the paradigm or exhausted its author's
weekend, and the record cannot say which. Replacing the author with an agent holds
the paradigm fixed and changes the cost of iteration --- the term that was never
varied.

\paragraph{Where to measure it.} The question needs a domain where many
independent authors pushed the symbolic paradigm as far as a real budget allowed,
with the results recorded and the budget the same for every method they might have
chosen. We use the complete archive of a collegiate game AI contest running since
\val{founded}: games authored for the contest itself, \val{nsub} submissions, and
\val{nmatch} adjudicated matches, written under deadline with standings at stake.
Contests of this kind carry a folk ordering of methods --- rules for beginners,
search for the intermediate, RL as the highest ceiling --- which is worth naming
because it is the ordering the field's own results are taken to endorse. The
archive records that this ordering gets the asymptotic ceilings right and still
fails to predict who wins: \val{pheur} of prize-winning programs contain no
learned component, and RL, present in \val{prl} of submissions, reached the finals
in \val{nrlfinal} of \val{editions} editions.

Two properties follow from provenance rather than design. Contamination resistance
is usually bought by authoring fresh tasks or collecting new ones continuously
\citep{livecodebench2024,workbuddybench2026}; here it comes from a venue that
predates the benchmark and was never built for evaluation --- which we still verify
rather than assume. And because thousands of people attempted the same games, the
archive carries a \emph{distribution} of human skill rather than a reference
solution, so a method can be placed at a percentile of a real field, as agents have
been placed against Kaggle competitors and informatics olympiad contestants
\citep{mlebench2024,olympiad2025percentile}.

\paragraph{Heuristic learning.} We remove the confound with heuristic learning
(HL): a coding agent maintains a heuristic program by reading the rules, its own
source, and replays of its losses, then rescores it against a frozen opponent pool.
The loop is deliberately narrow --- one revision per iteration, every version and
score kept --- because attribution is the property we are trying to preserve. Its
evaluator is a game engine and its opponents are frozen, placing it at the strong
end of the verification hierarchy along which demonstrated self-improvement tracks
\citep{selfimprovement2026survey}.

That this is not merely a matter of a stronger model matters, because the closest
existing head-to-head found the humans winning: across 38{,}304 matches,
human-written agents took the top five places and a frontier model handed the
winning student's program dropped it to tenth
\citep{vibecoding2025tournament}. Their agents were not one-shot --- most were
revised and all were debugged over many cycles --- but that loop optimized
compilability, and no agent was ever revised against the tournament it would be
judged in. The gap is therefore neither model capability nor a missing edit; it is
the absence of many revisions carrying an adversarial signal, which is the loop this
paper measures.

Two objections are conceded here and answered in full in
Section~\ref{sec:discussion}. The agent arrives with a pretraining corpus while RL
starts cold, and no denominator makes that symmetric; we report environment steps
as what is scarce in a new domain, and test what pretraining supplies
(Section~\ref{sec:experiments}) rather than asserting it. And the archive's RL
scarcity is confounded with its participants' cost structure --- students on a
deadline have limited compute, time, and RL expertise --- so what it records is
what wins under contest constraints, at one budget applied to whichever method a
contestant chose. Every comparison here is budget-conditional.

\subsection*{Contributions}

% Method, benchmark, comparison -- in that order, so the algorithm is the paper's
% subject and the benchmark is what it is measured on. The archive audit sits
% under the benchmark: it describes what the archive contains, and is not a
% result of comparing HL against RL.
\begin{itemize}[leftmargin=*]
  \item \textbf{Heuristic learning (HL): agent-maintained symbolic policies under
  adversarial pressure, instrumented so the ceiling can be located.} That an LLM
  can revise a program from execution feedback is established
  \citep{kuang2025generative}; what HL adds is the adversary and the discipline ---
  one revision per iteration, every version and score in an append-only log, so
  that a change in behaviour is attributable to the edit that caused it. Without
  that attribution a plateau can be observed but not diagnosed.

  \item \textbf{A benchmark drawn from a real competition rather than constructed
  for evaluation.} To the archive's \val{nsub} submissions and \val{nmatch}
  adjudicated matches --- written to win under deadline, not authored as tasks ---
  we add the apparatus a benchmark needs: classic games as structural controls,
  three opponent pools (training, frozen test, Elo anchor), a frozen evaluation
  set, and an Elo protocol comparable across games and editions. We also classify
  every submission by whether it contains a learned component, which is what makes
  the audit above a measurement rather than an impression.

  \item \textbf{A broad comparison of HL and RL, and a mechanism for where HL
  stops.} Against model-based RL, policy gradient, and the human field, over a
  range of budgets, reported in every denominator rather than the favourable one:
  HL reaches the \val{pct}th percentile of the field and beats the strongest
  recorded human on \val{nbeat} games, and is \val{effratio} more sample-efficient
  than model-based RL at matched Elo --- but that baseline overtakes beyond
  \val{crossover} environment steps and HL costs more \val{loseaxis}. Neither
  paradigm dominates; the budget and the game's structure decide. Improvement
  frequency decays as \val{decay} and plateaus near \val{plateau} revisions, which
  we trace to maintenance capacity on a pair of quantities: description length at
  the plateau converges to \val{maintcap} within \val{maintcapspread} across games
  of differing complexity, while the Elo reached there varies by
  \val{eloplatspread}. Equal maintainable size at unequal strength is what
  separates a limit of the agent from a limit of the game --- convergent length
  alone would not, being equally consistent with games that never differed. The
  loop transfers to \val{ntransfer} further benchmarks without retuning.
\end{itemize}

